\chapter{2007年重庆卷（理）}

\section{单选题}

\begin{problem}
若等差数列 \(\{a_{n}\}\) 的前 \(3\) 项和 \(S_{3} = 9\) 且 \(a_1 = 1\)，则 \(a_2\) 等于（\quad）

\choices
    {\(3\)}
    {\(4\)}
    {\(5\)}
    {\(6\)}
\end{problem}

\begin{problem}
命题“若 \(x^{2} < 1\)，则 \(-1 < x < 1\)”的逆否命题是（\quad）

\choices
    {若 \(x^{2} \geqslant 1\)，则 \(x \geqslant 1\) 或 \(x \leqslant -1\)}
    {若 \(-1 < x < 1\)，则 \(x^{2} < 1\)}
    {若 \(x > 1\) 或 \(x < -1\)，则 \(x^{2} > 1\)}
    {若 \(x \geqslant 1\) 或 \(x \leqslant -1\)，则 \(x^{2} \geqslant 1\)}
\end{problem}

\begin{problem}
若三个平面两两相交，且三条交线互相平行，则这三个平面把空间分成（\quad）

\choices
    {\(5\) 部分}
    {\(6\) 部分}
    {\(7\) 部分}
    {\(8\) 部分}
\end{problem}

\begin{problem}
若 \(\left(x + \frac{1}{x}\right)^n\) 展开式的二项式系数之和为 64，则展开式的常数项为（\quad）

\choices
    {\(10\)}
    {\(20\)}
    {\(30\)}
    {\(120\)}
\end{problem}

\begin{problem}
在 \(\triangle ABC\) 中，\(AB = \sqrt{3}\)，\(A = 45^{\circ}\)，\(C = 75^{\circ}\)，则 \(BC =\)（\quad）

\choices
    {\(3 - \sqrt{3}\)}
    {\(\sqrt{2}\)}
    {\(2\)}
    {\(3 + \sqrt{3}\)}
\end{problem}

\begin{problem}
从 \(5\) 张 \(100\text{ 元}\)，\(3\) 张 \(200\text{ 元}\)，\(2\) 张 \(300\text{ 元}\)的奥运预赛门票中任取 \(3\) 张，则所取 \(3\) 张中至少有 \(2\) 张价格相同的概率为（\quad）

\choices
    {\( \frac{1}{4} \)}
    {\( \frac{79}{120} \)}
    {\( \frac{3}{4} \)}
    {\( \frac{23}{24} \)}
\end{problem}

\begin{problem}
若 \(a\) 是 \(1 + 2b\) 与 \(1 - 2b\) 的等比中项，则 \(\frac{2ab}{|a| + 2|b|}\) 的最大值为（\quad）

\choices
    {\(\frac{2\sqrt{5}}{5}\)}
    {\(\frac{\sqrt{2}}{4}\)}
    {\(\frac{\sqrt{5}}{5}\)}
    {\(\frac{\sqrt{2}}{2}\)}
\end{problem}

\begin{problem}
设正数 \(a\)，\(b\) 满足 \(\displaystyle \lim\limits_{x \to 2} \left(x^{2} + ax - b\right) = 4\)，则 \(\displaystyle \lim\limits_{n \to \infty} \frac{a^{n+1} + ab^{n-1}}{a^{n-1} + 2b^n} =\)（\quad）

\choices
    {\(0\)}
    {\( \frac{1}{4} \)}
    {\( \frac{1}{2} \)}
    {\(1\)}
\end{problem}

\begin{problem}
已知定义域为 \(\R\) 的函数 \(f(x)\) 在 \((8, +\infty)\) 上为减函数，且函数 \(y = f(x + 8)\) 为偶函数，则（\quad）

\choices
    {\(f(6) > f(7)\)}
    {\(f(6) > f(9)\)}
    {\(f(7) > f(9)\)}
    {\(f(7) > f(10)\)}
\end{problem}

\begin{problem}
如图，在四边形 \(ABCD\) 中，\(\left|\overrightarrow{AB}\right| + \left|\overrightarrow{BD}\right| + \left|\overrightarrow{DC}\right| = 4\)，\(\left|\overrightarrow{AB}\right| \cdot \left|\overrightarrow{BD}\right| + \left|\overrightarrow{BD}\right| \cdot \left|\overrightarrow{DC}\right| = 4\)，\(\overrightarrow{AB} \cdot \overrightarrow{BD} = \overrightarrow{BD} \cdot \overrightarrow{DC} = 0\)，则 \(\left(\overrightarrow{AB} + \overrightarrow{DC}\right) \cdot \overrightarrow{AC}\) 的值为（\quad）

\bitmapfigure[max width=0.34\linewidth]{img/2007/chongqing_science/q10_fig1.png}

\choices
    {\(2\)}
    {\(2 \sqrt{2}\)}
    {\(4\)}
    {\(4 \sqrt{2}\)}
\end{problem}

\section{填空题}

\begin{problem}
复数 \(\frac{2\i}{2 + \i^{3}}\) 的虚部为\fillinblank{}.
\end{problem}

\begin{problem}
已知 \(x\)，\(y\) 满足 \(\begin{cases} x - y \leqslant 1, \\ 2x + y \leqslant 4, \\ x \geqslant 1, \end{cases}\) 则函数 \(z = x + 3y\) 的最大值是\fillinblank{}.
\end{problem}

\begin{problem}
若函数 \( f(x) = \sqrt{2^{x^{2} + 2ax - a} - 1} \) 的定义域为 \( \R \)，则 \( a \) 的取值范围为\fillinblank{}.
\end{problem}

\begin{problem}
设 \( \{a_{n}\} \) 为公比 \(q > 1\) 的等比数列，若 \( a_{2004} \) 和 \( a_{2006} \) 是方程 \( 4x^{2} - 8x + 3 = 0 \) 的两根，则 \( a_{2006} + a_{2007} = \)\fillinblank{}.
\end{problem}

\begin{problem}
某校要求每位学生从 \(7\) 门课程中选修 \(4\) 门，其中甲，乙两门课程不能都选，则不同的选课方案有\fillinblank{} 种. （用数字作答）
\end{problem}

\begin{problem}
过双曲线 \( x^{2}-y^{2}=4 \) 的右焦点 \(F\) 作倾斜角为 \( 105^{\circ} \) 的直线，交双曲线于 \(P\)，\(Q\) 两点，则 \( |FP| \cdot |FQ| \) 的值为\fillinblank{}.
\end{problem}

\section{解答题}

\begin{problem}
设 \( f(x) = 6\cos^2 x - \sqrt{3}\sin 2x \).

\begin{enumerate}
\item 求 \( f(x) \) 的最大值及最小正周期；
\item 若锐角 \(\alpha\) 满足 \(f(\alpha) = 3 - 2\sqrt{3}\)，求 \(\tan \frac{4}{5} \alpha\) 的值.
\end{enumerate}
\end{problem}

\begin{problem}
某单位有三辆汽车参加某种事故保险. 单位年初向保险公司缴纳每辆 \(900\text{ 元}\) 的保险金. 对在一年内发生此种事故的每辆汽车，单位可获 \(9000\text{ 元}\) 的赔偿（假设每辆车最多只赔偿一次）. 设这三辆车在一年内发生此种事故的概率分别为 \(\frac{1}{9}\)，\(\frac{1}{10}\)，\(\frac{1}{11}\)，且各车是否发生事故相互独立. 求一年内该单位在此保险中：

\begin{enumerate}
\item 获赔的概率；
\item 获赔金额 \( \xi \) 的分布列与期望.
\end{enumerate}
\end{problem}

\begin{problem}
如图，在直三棱柱 \(ABC - A_{1}B_{1}C_{1}\) 中，\(AA_{1} = 2\)，\(AB = 1\)，\(\angle ABC = 90^{\circ}\)；点 \(D\)，\(E\) 分别在 \(BB_{1}\)，\(A_{1}D\) 上，且 \(B_{1}E \perp A_{1}D\)，四棱锥 \(C - ABDA_{1}\) 与直三棱柱的体积之比为 \(3:5\).

\begin{enumerate}
\item 求异面直线 \(DE\) 与 \(B_{1}C_{1}\) 的距离；
\item 若 \(BC = \sqrt{2}\)，求二面角 \(A_{1} - DC_{1} - B_{1}\) 的平面角的正切值.
\end{enumerate}
\bitmapfigure[max width=0.34\linewidth]{img/2007/chongqing_science/q19_fig1.png}
\end{problem}

\begin{problem}
已知函数 \( f(x) = ax^{4}\ln x + bx^{4} - c \)（\(x > 0\)）在 \( x = 1 \) 处取得极值 \(-3 - c\)，其中 \(a\)，\(b\)，\(c\) 为常数.

\begin{enumerate}
\item 试确定 \(a\)，\(b\) 的值；
\item 讨论函数 \( f(x) \) 的单调区间；
\item 若对任意 \(x > 0\)，不等式 \(f(x) \geqslant -2c^{2}\) 恒成立，求 \(c\) 的取值范围.
\end{enumerate}
\end{problem}

\begin{problem}
已知各项均为正数的数列 \(\{a_{n}\}\) 的前 \(n\) 项和 \(S_{n}\) 满足 \(S_{1} > 1\)，且 \(6S_{n} = (a_{n} + 1)(a_{n} + 2)\)，\(n \in \mathbf{N}_{+}\).

\begin{enumerate}
\item 求 \( \{a_{n}\} \) 的通项公式；
\item 设数列 \( \{b_{n}\} \) 满足 \( a_{n}(2^{b_{n}}-1)=1 \)，并记 \( T_{n} \) 为 \( \{b_{n}\} \) 的前 \(n\) 项和，求证： \(3T_{n}+1>\log_{2}(a_{n}+3)\)，\(n\in\mathbf{N}_{+}\).
\end{enumerate}
\end{problem}

\begin{problem}
如图，中心在原点 \(O\) 的椭圆的右焦点为 \(F(3,0)\)，右准线 \(l\) 的方程为：\(x = 12\).

\begin{enumerate}
\item 求椭圆的方程；
\item 在椭圆上任取三个不同点 \(P_{1}\)，\(P_{2}\)，\(P_{3}\)，使 \(\angle P_{1}FP_{2} = \angle P_{2}FP_{3} = \angle P_{3}FP_{1}\)，证明 \(\frac{1}{|FP_1|} + \frac{1}{|FP_2|} + \frac{1}{|FP_3|}\) 为定值，并求此定值.
\end{enumerate}
\bitmapfigure[max width=0.38\linewidth]{img/2007/chongqing_science/q22_fig1.png}
\end{problem}
